\documentclass[%
 amsmath,amssymb,
 aps,
 jor,
 sor,
 nofootinbib,
12pt,]{revtex4-2}
\usepackage{amsmath}
\usepackage{hyperref}
\usepackage{etoolbox}
\usepackage{titletoc}
\usepackage{ragged2e}

\titlecontents{section}[0pt]
{\bfseries}
{\contentslabel{2em}}
{}
{\titlerule*[0.5pc]{}\contentspage}

\usepackage{graphicx}
\usepackage{dcolumn}
\usepackage{bm}
\usepackage{float}
\usepackage{mathtools}
\usepackage{hyperref}
\hypersetup{
     colorlinks=true,
     linkcolor=blue,
     filecolor=blue,
     citecolor=blue,
     urlcolor=blue,
}
\usepackage[mathlines]{lineno}

\begin{document}
\preprint{APS/123-QED}
\title{Are LSTM Uncertainty Estimates Calibrated?\\A Post-Hoc Diagnostic Study of Probabilistic WTI Crude Oil Forecasts}
\author{Peter Dunmain, SN: XXXXXXX, COMP0047}
\maketitle
\renewenvironment{abstract}{
    \begin{center}
    \textbf{\large Abstract}
    \end{center}
    \par
}{}
\begin{center}
    \textbf{ABSTRACT}
\end{center}
\begin{quote}
The group project claimed that the LSTM's learned $\sigma_t$ is ``calibrated and regime-aware,'' but provided no formal statistical evidence.
This study applies a battery of goodness-of-fit tests to the standardised residuals of two LSTM variants and a GARCH(1,1) benchmark, finding that the pure-NLL LSTM produces partially calibrated uncertainty estimates while the composite decision-aware loss materially distorts $\sigma_t$.
A Diebold--Mariano test confirms that the NLL LSTM significantly outperforms the composite model ($p < 0.001$) but does not significantly outperform GARCH ($p = 0.67$), and a Student-$t$ extension consistently improves calibration for all models.
\end{quote}

\tableofcontents
\newpage

%% =========================================================================
\section{Introduction}
\label{sec:intro}
%% =========================================================================

Forecasting next-day returns in commodity markets is a task characterised by extremely low signal-to-noise ratios.
When models are trained under standard loss functions such as mean squared error or Gaussian negative log-likelihood (NLL), the optimal prediction is typically close to zero --- the unconditional mean of the return distribution.
This phenomenon, which the group project referred to as the ``flat-line problem,'' is not a failure of modelling but a consequence of the signal structure: the conditional mean of daily log returns is very close to zero, and any aggressive departure risks large quadratic penalties.

The group project addressed this problem by designing a composite three-component loss function that combines Gaussian NLL with a directional penalty and a variance-collapse penalty.
Under this composite loss, the LSTM produced non-trivial directional predictions (52.3\% accuracy versus 47.1\% under pure NLL) and a ten-fold increase in prediction spread.
The group then claimed that the LSTM's learned $\sigma_t$ is ``calibrated and regime-aware,'' citing two pieces of informal evidence: the unconditional average of $\sigma_t$ ($\approx 0.025$) is close to the empirical return standard deviation ($\approx 0.023$), and the uncertainty bands visually widen during volatile periods.

Neither of these observations constitutes a formal calibration test.
Matching unconditional moments is a necessary but not sufficient condition for probabilistic calibration; visual band-widening is not quantitative.
Moreover, the composite loss is explicitly not a proper scoring rule in the sense of \citet{gneiting2007strictly}.
To see why, note that a scoring rule $S(F, r)$ is \emph{proper} if its expectation is minimised when the forecaster reports the true predictive distribution $G$: $\mathbb{E}_{r \sim G}[S(G, r)] \leq \mathbb{E}_{r \sim G}[S(F, r)]$ for all $F$.
The Gaussian NLL satisfies this property because it is the logarithmic score for the Gaussian family.
But the composite loss adds terms --- $\mathcal{L}_{\text{dir}}$ and $\mathcal{L}_{\text{var}}$ --- that depend on $\mu_t$ through non-probabilistic objectives.
The expectation of these terms under the true distribution is not minimised at $F = G$, so the composite loss cannot be proper: the model is incentivised to report a distribution different from its best estimate of the truth, and $\sigma_t$ absorbs the resulting distortion.

Calibration of neural network uncertainty estimates is a known challenge in the broader deep learning literature \citep{lakshminarayanan2017simple,kuleshov2018accurate}, and the specific question of whether heteroscedastic regression networks produce calibrated variance estimates has received less attention in the financial forecasting context.
The group additionally proposed using $\sigma_t$ for position sizing in a simulated trading strategy as future work.
Before $\sigma_t$ can be trusted for downstream risk management, its calibration must be established --- or its miscalibration must be understood.

This report addresses the following research question, which I consider the most important open question from the group project:

\medskip
\noindent\textbf{Research question.} \emph{Are the LSTM's learned uncertainty estimates actually calibrated, and does calibration survive the transition from the pure Gaussian NLL objective to the composite decision-aware loss?}
\medskip

Three hypotheses guide my analysis:

\begin{itemize}
    \item[\textbf{H1}:] Under pure NLL, the standardised residuals $z_t = (r_t - \mu_t) / \sigma_t$ are approximately consistent with the assumed $\mathcal{N}(0,1)$ predictive law.
    \item[\textbf{H2}:] Calibration deteriorates under the composite loss because the competing objectives distort the probabilistic interpretation of $\sigma_t$.
    \item[\textbf{H3}:] A Student-$t$ predictive distribution improves calibration under both objectives because WTI returns exhibit excess kurtosis even after standardisation.
\end{itemize}

The validation toolkit draws directly from the course lecture material, primarily Chapter~18 (Assessing the Goodness of Models) and Chapter~14 (Parametric Estimation).
I use Q-Q plots (Section~18.3), the Kolmogorov--Smirnov test (Section~18.3.3), the Anderson--Darling test (Section~18.3.4), likelihood-based model comparison (Section~18.6), the Akaike Information Criterion (Section~18.6), and the Ljung--Box test from the stochastic processes framework of Chapter~9.
I supplement these with two tests not covered in lectures: the Diebold--Mariano test \citep{diebold1995comparing} for formal comparison of predictive densities, and the Christoffersen conditional coverage test \citep{christoffersen1998evaluating} for detecting clustered violations.
The Student-$t$ maximum likelihood extension connects to the parametric estimation framework of Chapter~14.

%% =========================================================================
\section{Data Analysis}
\label{sec:data}
%% =========================================================================

\subsection{Dataset and target variable}

The dataset contains 1,895 daily observations of WTI crude oil and related features spanning 31 January 2017 to 3 March 2026.
The target variable is the next-day WTI log return:
\begin{equation}
    r_t = \log\!\left(\frac{P_{t+1}}{P_t}\right),
    \label{eq:target}
\end{equation}
where $P_t$ is the WTI spot price on day $t$.
The feature set used for the probabilistic LSTM models comprises 51 inputs from the PCA-compressed dataset, spanning five groups: market prices and log returns (WTI, Brent, S\&P~500, VIX, DXY), lagged WTI returns (lags 1--10), rolling volatility (5-, 10-, 20-, 60-day windows), EIA inventory statistics, and three PCA components summarising 28 Federal Reserve text sentiment features.
All features are strictly point-in-time; a leakage audit confirmed the absence of look-ahead bias.

\subsection{Train/validation/test split}

The data is split chronologically into three non-overlapping periods, as shown in Table~\ref{tab:splits}.
The \texttt{RobustScaler} for both features and targets is fit exclusively on the training set and applied to validation and test data, ensuring no information leakage.

\begin{table}[H]
\centering
\caption{Walk-forward data split.}
\label{tab:splits}
\begin{ruledtabular}
\begin{tabular}{lccc}
\textbf{Split} & \textbf{Period} & \textbf{Observations} & \textbf{Share} \\
\hline
Train      & 31 Jan 2017 -- 30 Dec 2022 & 1,235 & 65\% \\
Validation & 3 Jan 2023 -- 28 Jun 2024  & 312   & 16\% \\
Test       & 1 Jul 2024 -- 3 Mar 2026   & 348   & 18\% \\
\end{tabular}
\end{ruledtabular}
\end{table}

\subsection{Test-set return distribution}

The test-set return distribution exhibits features characteristic of financial time series.
Table~\ref{tab:return_stats} reports the summary statistics.
The negative skewness ($-0.77$) and excess kurtosis ($4.80$) confirm that the Gaussian distribution is a poor unconditional model for these returns, consistent with the classical stylised facts of asset returns \citep{cont2001empirical}.
Fitting a Student-$t$ distribution by maximum likelihood yields $\hat{\nu} = 5.73$, $\hat{\mu} = -0.00004$, and $\hat{\sigma} = 0.0183$, indicating fat tails that persist even in the most recent test period.

\begin{table}[H]
\centering
\caption{Summary statistics for test-set returns ($n = 348$).}
\label{tab:return_stats}
\begin{ruledtabular}
\begin{tabular}{lcccc}
& \textbf{Mean} & \textbf{Std} & \textbf{Skewness} & \textbf{Excess kurtosis} \\
\hline
$r_t$ & $-0.0004$ & $0.0230$ & $-0.77$ & $4.80$ \\
\end{tabular}
\end{ruledtabular}
\end{table}

\subsection{Group project models}

I do not train any new models in this analysis.
It uses the test-set outputs of two probabilistic LSTM variants from the group project, each producing a conditional mean $\mu_t$ and conditional standard deviation $\sigma_t$ for every test day:

\begin{enumerate}
    \item \textbf{NLL LSTM} --- trained under pure Gaussian NLL using a stacked LSTM architecture \citep{hochreiter1997long} with a dual-headed output layer producing $(\mu_t, \log\sigma_t^2)$ following the heteroscedastic regression framework of \citet{nix1994estimating}.
    Architecture: 3-layer stacked LSTM, hidden size 128, sequence length 3, no dropout, learning rate $10^{-3}$.
    The loss function is
    \begin{equation}
        \mathcal{L}_{\text{NLL}} = \frac{1}{2}\left[\log\sigma_t^2 + \frac{(r_t - \mu_t)^2}{\sigma_t^2}\right].
        \label{eq:nll}
    \end{equation}
    \item \textbf{Composite LSTM} --- trained under $\mathcal{L} = \mathcal{L}_{\text{NLL}} + \beta \mathcal{L}_{\text{dir}} + \gamma \mathcal{L}_{\text{var}}$, where
    \begin{equation}
        \mathcal{L}_{\text{dir}} = 1 - \overline{\sigma\!\left(\tau \cdot \mu_t \cdot r_t\right)}, \qquad
        \mathcal{L}_{\text{var}} = \frac{\gamma}{\operatorname{Var}(\mu_t) + \epsilon},
        \label{eq:composite}
    \end{equation}
    with $\beta = 1.0$ and $\gamma = 0.05$.
    Architecture: 2-layer stacked LSTM, hidden size 128, sequence length 5, dropout 0.3, learning rate $5 \times 10^{-4}$.
\end{enumerate}

In addition, a \textbf{GARCH(1,1)} model is fit to the training returns as a classical volatility benchmark.
GARCH is the standard tool for modelling conditional heteroscedasticity in financial returns, building on Engle's original ARCH framework \citep{engle1982autoregressive,bollerslev1986generalized}, and provides a natural reference point.
The group project's slides explicitly connect the LSTM's heteroscedastic output to GARCH-style modelling, making this comparison directly relevant.

Figure~\ref{fig:eda} presents the test-set return distribution alongside fitted Gaussian and Student-$t$ density overlays, a Q-Q plot, a P-P plot, and the rolling 60-day realised volatility.
The Q-Q plot reveals clear tail departures from the Gaussian model, particularly in the left tail.
The rolling volatility panel shows distinct volatility regimes during the test period, providing the visual context for assessing whether the models' $\sigma_t$ tracks these regimes.

\begin{figure}[H]
    \centering
    \includegraphics[width=\textwidth]{../notebooks/daily/return_prediction/figures/figure1_return_distribution_eda.pdf}
    \caption{Test-set return distribution.
    Top left: histogram with Gaussian and Student-$t$ MLE overlays.
    Top right: Q-Q plot of standardised $r_t$ against $\mathcal{N}(0,1)$.
    Bottom left: P-P plot of $r_t$ against the fitted Gaussian.
    Bottom right: rolling 60-day realised volatility.}
    \label{fig:eda}
\end{figure}

%% =========================================================================
\section{Methodology}
\label{sec:methodology}
%% =========================================================================

\subsection{Probabilistic calibration}

A probabilistic forecaster that outputs $(\mu_t, \sigma_t)$ is \emph{calibrated} if, for every nominal coverage level $\alpha \in (0,1)$, the empirical fraction of observations falling within the $\alpha$-prediction interval equals $\alpha$:
\begin{equation}
    \frac{1}{n}\sum_{t=1}^{n} \mathbf{1}\!\left[r_t \in \left(\mu_t - q_{\alpha}\,\sigma_t,\;\mu_t + q_{\alpha}\,\sigma_t\right)\right] \approx \alpha,
    \label{eq:coverage}
\end{equation}
where $q_{\alpha}$ is the $(1+\alpha)/2$ quantile of the assumed predictive distribution.

An equivalent characterisation is provided by the probability integral transform (PIT).
If the predictive distribution $F_t(r) = \Phi\!\left((r - \mu_t)/\sigma_t\right)$ is correctly specified, the PIT values
\begin{equation}
    u_t = \Phi\!\left(\frac{r_t - \mu_t}{\sigma_t}\right)
    \label{eq:pit}
\end{equation}
should be independently and identically distributed as $\text{Uniform}[0,1]$ \citep{diebold1998evaluating}.
Equivalently, the standardised residuals
\begin{equation}
    z_t = \frac{r_t - \mu_t}{\sigma_t}
    \label{eq:z}
\end{equation}
should be i.i.d.\ $\mathcal{N}(0,1)$ under a correctly specified Gaussian model.
Any departure from this reference law reveals a specific calibration failure: fat tails in $z_t$ indicate that $\sigma_t$ is too small on extreme days; serial correlation in $|z_t|$ indicates that $\sigma_t$ does not properly track the volatility process; empirical coverage below nominal indicates overconfidence.

\subsection{Statistical tests}

The calibration battery applies the following tests, each drawn from the course material.

\subsubsection{Kolmogorov--Smirnov test (Chapter 18, Section 18.3.3)}

The KS test computes the supremum distance between the empirical CDF $\hat{F}(x)$ and the hypothesised CDF $\tilde{F}(x)$:
\begin{equation}
    D_n = \sup_{x}\left|\hat{F}(x) - \tilde{F}(x)\right|.
    \label{eq:ks}
\end{equation}
Under the null hypothesis that the sample comes from $\tilde{F}$, the statistic $\sqrt{n}\,D_n$ converges to the Kolmogorov distribution.
The KS test is sensitive to shifts in location and scale but less sensitive to tail departures because the CDF is close to 0 or 1 in the tails.

\subsubsection{Anderson--Darling test (Chapter 18, Section 18.3.4)}

The AD test addresses the tail-sensitivity weakness of the KS test by introducing a weighting function that emphasises the tails:
\begin{equation}
    A^2 = -n - \frac{1}{n}\sum_{s=1}^{n}(2s-1)\left[\ln\tilde{F}(x_{(s)}) + \ln\!\left(1 - \tilde{F}(x_{(n+1-s)})\right)\right],
    \label{eq:ad}
\end{equation}
where $x_{(1)} \leq \cdots \leq x_{(n)}$ are the order statistics.
For the Student-$t$ tests, where critical values are not tabulated, the $p$-value is computed by a parametric Monte Carlo bootstrap: 9,999 samples of size $n$ are drawn from the fitted Student-$t$ distribution, the AD statistic is recomputed for each sample after re-fitting the degrees-of-freedom parameter, and the $p$-value is the fraction of simulated statistics exceeding the observed statistic.

\subsubsection{Ljung--Box test (Chapter 9)}

The Ljung--Box test \citep{ljung1978measure} checks for residual autocorrelation at lag $K$:
\begin{equation}
    Q(K) = n(n+2)\sum_{k=1}^{K}\frac{\hat{\rho}_k^2}{n-k},
    \label{eq:lb}
\end{equation}
where $\hat{\rho}_k$ is the sample autocorrelation at lag $k$.
Under the null of no autocorrelation, $Q(K) \sim \chi^2(K)$.
The test is applied both to $z_t$ (testing for residual mean structure) and to $|z_t|$ (testing for residual volatility clustering).
Significant autocorrelation in $|z_t|$ indicates that $\sigma_t$ fails to capture the full volatility dynamics.

\subsubsection{Interval coverage test}

The empirical coverage at nominal levels $\alpha \in \{0.50, 0.80, 0.90, 0.95, 0.99\}$ is computed via Equation~\eqref{eq:coverage}.
The deviation from nominal coverage provides a direct, interpretable measure of calibration quality \citep{christoffersen1998evaluating}.

\subsubsection{Out-of-sample negative log-likelihood (Chapter 18, Section 18.6)}

The out-of-sample NLL per observation under the Gaussian predictive law is
\begin{equation}
    \text{NLL}_{\text{Gauss}} = \frac{1}{n}\sum_{t=1}^{n}\left[\frac{1}{2}\log(2\pi) + \log\sigma_t + \frac{1}{2}\!\left(\frac{r_t - \mu_t}{\sigma_t}\right)^{\!2}\right],
    \label{eq:gauss_nll}
\end{equation}
where $\phi(\cdot)$ is the standard normal density.
Lower values indicate a better probabilistic fit.
The NLL provides a single scalar summary of probabilistic forecast quality and connects directly to the cross-entropy framework of Section~18.5.2.

\subsubsection{Christoffersen conditional coverage test}

The unconditional coverage test checks only whether the \emph{fraction} of violations equals the nominal rate.
\citet{christoffersen1998evaluating} showed that a stronger requirement is that violations be \emph{independent}: if today's violation predicts tomorrow's, the model's $\sigma_t$ is not adapting fast enough to the current regime.
The test models the violation sequence as a first-order Markov chain with transition probabilities $\pi_{01}$ (violation following no violation) and $\pi_{11}$ (violation following violation).
Under the null of independence, $\pi_{01} = \pi_{11}$, and the likelihood ratio statistic
\begin{equation}
    \text{LR}_{\text{ind}} = 2\left[\ell_1(\hat{\pi}_{01}, \hat{\pi}_{11}) - \ell_0(\hat{\pi})\right] \sim \chi^2(1),
    \label{eq:cc}
\end{equation}
where $\ell_1$ is the unrestricted (Markov) log-likelihood and $\ell_0$ is the restricted (independence) log-likelihood.

\subsubsection{Diebold--Mariano test for predictive density comparison}

The Diebold--Mariano (DM) test \citep{diebold1995comparing} provides a formal statistical comparison of predictive densities.
For two models with log-scores $s_t^{(1)}$ and $s_t^{(2)}$ at each time $t$, the test statistic is based on the loss differential $d_t = s_t^{(1)} - s_t^{(2)}$:
\begin{equation}
    \text{DM} = \frac{\bar{d}}{\hat{\sigma}_d / \sqrt{n}},
    \label{eq:dm}
\end{equation}
where $\bar{d}$ is the sample mean of $d_t$ and $\hat{\sigma}_d$ is a Newey--West consistent variance estimator that accounts for possible serial correlation in $d_t$.
Under the null hypothesis of equal predictive ability, DM is asymptotically standard normal.
Unlike the raw NLL comparison, the DM test answers whether the observed NLL differences are statistically significant or could arise from sampling variability.

\subsubsection{Akaike Information Criterion (Chapter 18, Section 18.6)}

The AIC penalises the log-likelihood by the number of free distributional parameters $k$:
\begin{equation}
    \text{AIC} = 2k - 2\ell,
    \label{eq:aic}
\end{equation}
where $\ell$ is the total log-likelihood.
For the Gaussian specification, $k = 0$ (the model's $\mu_t$ and $\sigma_t$ are given as inputs, not free parameters of the distributional fit).
For the Student-$t$ specification, $k = 1$ (the degrees-of-freedom parameter $\nu$ is estimated from the standardised residuals).

\subsection{GARCH(1,1) benchmark}

The GARCH(1,1) model \citep{bollerslev1986generalized} specifies the conditional variance as
\begin{equation}
    \sigma_t^2 = \omega + \alpha\,r_{t-1}^2 + \beta\,\sigma_{t-1}^2,
    \label{eq:garch}
\end{equation}
where $\omega > 0$, $\alpha \geq 0$, $\beta \geq 0$, and $\alpha + \beta < 1$ ensures stationarity.
The model is fit by maximum likelihood on the training returns with a zero-mean specification, reflecting the fact that the unconditional mean of daily log returns is negligibly small.
Out-of-sample forecasts for the test period are produced by recursively applying Equation~\eqref{eq:garch}, bridging through the validation period to preserve the volatility state.

\subsection{Conditional calibration}

Unconditional calibration tests may mask regime-dependent miscalibration.
If a model is overconfident during volatile periods and overdispersed during calm periods, the aggregate statistics may average to acceptable values.
To detect such behaviour, the test set is split into high-volatility and low-volatility subsets using the median of the rolling 60-day realised volatility as the threshold, following the conditional calibration framework of \citet{kuleshov2018accurate}.
The calibration battery is then applied separately to each subset.

\subsection{Student-$t$ extension (Chapter 14)}

After documenting Gaussian miscalibration, the analysis relaxes the distributional assumption.
For each model's standardised residuals $z_t$, the degrees-of-freedom parameter $\hat{\nu}$ is estimated by maximum likelihood, maximising
\begin{equation}
    \ell(\nu) = \sum_{t=1}^{n}\log f_{\nu}(z_t),
    \label{eq:t_mle}
\end{equation}
where $f_{\nu}$ is the standard Student-$t$ density with $\nu$ degrees of freedom, zero location, and unit scale.
This connects directly to the parametric estimation framework of Chapter~14.
The entire calibration battery is then rerun using the Student-$t$ predictive law in place of the Gaussian.

\subsection{Limitations}

The test period (July 2024 -- March 2026) was already used for final evaluation in the group project.
This individual analysis treats it as a post-hoc diagnostic sample rather than a fresh holdout.
The calibration results are therefore descriptive of these specific model outputs rather than a fully independent validation.

%% =========================================================================
\section{Results}
\label{sec:results}
%% =========================================================================

\subsection{Standardised residual diagnostics}

Table~\ref{tab:z_stats} reports the summary statistics of the standardised residuals for all three models.
Under perfect Gaussian calibration, the residuals should have mean 0, standard deviation 1, skewness 0, and excess kurtosis 0.

\begin{table}[H]
\centering
\caption{Summary statistics of standardised residuals $z_t$.}
\label{tab:z_stats}
\begin{ruledtabular}
\begin{tabular}{lcccc}
\textbf{Model} & \textbf{Mean} & \textbf{Std} & \textbf{Skewness} & \textbf{Excess kurtosis} \\
\hline
NLL LSTM       & $-0.070$ & $0.938$ & $-0.89$ & $5.01$ \\
Composite LSTM & $-0.013$ & $0.904$ & $-0.18$ & $3.00$ \\
GARCH(1,1)     & $-0.014$ & $1.039$ & $-1.40$ & $11.64$ \\
\end{tabular}
\end{ruledtabular}
\end{table}

Several observations follow.
The NLL LSTM residuals have a slight negative mean bias ($-0.070$) and a standard deviation below unity ($0.938$), indicating mild overdispersion of $\sigma_t$ relative to the actual forecast errors.
The excess kurtosis of $5.01$ confirms persistent fat tails even after standardisation, meaning the Gaussian distributional assumption is violated.
The composite LSTM residuals have a standard deviation further below unity ($0.904$), indicating stronger overdispersion, and an excess kurtosis of $3.00$, still well above the Gaussian reference of zero.
The GARCH residuals exhibit the most extreme departure from normality: excess kurtosis of $11.64$ and skewness of $-1.40$, indicating that while GARCH tracks the volatility level, it does not adequately capture the most extreme moves.

Figure~\ref{fig:residuals} displays the time series of $z_t$, the histogram against the $\mathcal{N}(0,1)$ reference, and the predicted $\sigma_t$ against the 60-day rolling realised volatility for both LSTM models.

\begin{figure}[H]
    \centering
    \includegraphics[width=\textwidth]{../notebooks/daily/return_prediction/figures/figure2_standardised_residuals.pdf}
    \caption{Standardised residuals and predicted volatility for the NLL LSTM (top row) and composite LSTM (bottom row).
    Left: $z_t$ over time with $\pm 1.96$ bands.
    Centre: histogram against $\mathcal{N}(0,1)$.
    Right: predicted $\sigma_t$ versus 60-day rolling realised volatility.}
    \label{fig:residuals}
\end{figure}

\subsection{Autocorrelation diagnostics}

Figure~\ref{fig:acf} presents the ACF of $z_t$ and $|z_t|$ for both LSTM models, and Table~\ref{tab:lb} reports the Ljung--Box test results at lag 10.

\begin{figure}[H]
    \centering
    \includegraphics[width=\textwidth]{../notebooks/daily/return_prediction/figures/figure3_acf_grid.pdf}
    \caption{Autocorrelation functions of $z_t$ (left column) and $|z_t|$ (right column) for the NLL LSTM (top) and composite LSTM (bottom).
    Dashed lines indicate the 95\% confidence band under the white-noise null.}
    \label{fig:acf}
\end{figure}

\begin{table}[H]
\centering
\caption{Ljung--Box test at lag 10 for the standardised residuals and their absolute values.}
\label{tab:lb}
\begin{ruledtabular}
\begin{tabular}{llcc}
\textbf{Model} & \textbf{Series} & $Q(10)$ & $p$\textbf{-value} \\
\hline
NLL LSTM       & $z_t$    & 9.52  & 0.48 \\
NLL LSTM       & $|z_t|$  & 7.24  & 0.70 \\
Composite LSTM & $z_t$    & 16.35 & 0.09 \\
Composite LSTM & $|z_t|$  & 5.98  & 0.82 \\
GARCH(1,1)     & $z_t$    & 8.14  & 0.61 \\
GARCH(1,1)     & $|z_t|$  & 7.02  & 0.72 \\
\end{tabular}
\end{ruledtabular}
\end{table}

None of the $|z_t|$ series show significant autocorrelation at the 5\% level, suggesting that all three models capture the broad volatility dynamics without leaving systematic volatility clustering in the residuals.
However, the composite LSTM's $z_t$ series shows borderline significance ($p = 0.09$), indicating some residual serial dependence in the mean predictions.
A sensitivity check across lag choices (5, 10, 15, 20) confirms that this borderline result is robust: the composite LSTM's $z_t$ $p$-value ranges from $0.057$ to $0.173$, consistently near the significance threshold, while all other model-series combinations remain clearly non-significant.

\subsection{Formal distributional tests}

Table~\ref{tab:ks_ad} reports the KS and AD test results for the standardised residuals against $\mathcal{N}(0,1)$.

\begin{table}[H]
\centering
\caption{KS and Anderson--Darling tests of $z_t$ against $\mathcal{N}(0,1)$.}
\label{tab:ks_ad}
\begin{ruledtabular}
\begin{tabular}{lcccc}
\textbf{Model} & $D_n$ \textbf{(KS)} & $p$\textbf{-value (KS)} & $A^2$ \textbf{(AD)} & \textbf{AD at 5\%} \\
\hline
NLL LSTM       & 0.083 & 0.015 & 1.456 & Rejected \\
Composite LSTM & 0.063 & 0.123 & 1.160 & Rejected \\
GARCH(1,1)     & 0.043 & 0.534 & 2.081 & Rejected \\
\end{tabular}
\end{ruledtabular}
\end{table}

The results reveal an informative divergence between the two tests.
The KS test rejects the Gaussian null for the NLL LSTM ($p = 0.015$) but fails to reject for the composite LSTM ($p = 0.123$) and GARCH ($p = 0.534$).
However, the AD test --- which weights the tails more heavily --- rejects for all three models.
This pattern is consistent with the known limitation of the KS test discussed in the lectures (Section~18.3.3): it is insensitive to tail departures because the CDF is close to 0 or 1 in the tails.
The AD test, by contrast, is designed precisely to detect tail misspecification, and its rejection across all models confirms that the Gaussian distributional assumption is inadequate for WTI return residuals regardless of the model producing them.

The fact that the KS test \emph{fails to reject} for GARCH despite GARCH having the highest excess kurtosis (11.64) illustrates this limitation clearly.

\subsection{PIT histograms}

The probability integral transform (PIT) provides a complementary visual diagnostic \citep{diebold1998evaluating}.
If the predictive CDF is correctly specified, the PIT values $u_t = \Phi(z_t)$ should be uniformly distributed on $[0,1]$.
Departures from uniformity reveal specific calibration failures: a U-shape indicates underdispersion (intervals too narrow), a hump shape indicates overdispersion (intervals too wide), and asymmetry indicates location bias.

Figure~\ref{fig:pit} presents the PIT histograms for all three models under the Gaussian predictive law.
All three models exhibit a hump-shaped PIT histogram, confirming the overdispersion identified in the coverage analysis.
The NLL LSTM and composite LSTM both show mild asymmetry in the left tail, consistent with the negative skewness of their standardised residuals.
A KS test of the PIT values against $\text{Uniform}[0,1]$ rejects for the NLL LSTM ($p = 0.015$) but fails to reject for the composite LSTM ($p = 0.123$) and GARCH ($p = 0.534$), mirroring the distributional test results.

\begin{figure}[H]
    \centering
    \includegraphics[width=\textwidth]{../notebooks/daily/return_prediction/figures/figure3b_pit_histograms.pdf}
    \caption{PIT histograms under the Gaussian predictive law.
    The dashed horizontal line indicates the uniform reference density.
    A perfectly calibrated model produces a flat histogram.}
    \label{fig:pit}
\end{figure}

\subsection{Q-Q diagnostics}

Figure~\ref{fig:qq} presents the Q-Q plots of $z_t$ against $\mathcal{N}(0,1)$ (top row) and against the fitted Student-$t$ distribution (bottom row), with 95\% pointwise confidence bands.

\begin{figure}[H]
    \centering
    \includegraphics[width=\textwidth]{../notebooks/daily/return_prediction/figures/figure4_qq_grid.pdf}
    \caption{Q-Q plots.
    Top row: $z_t$ vs $\mathcal{N}(0,1)$ for the NLL LSTM (left) and composite LSTM (right).
    Bottom row: $z_t$ vs the fitted Student-$t$ distribution.
    Grey bands indicate 95\% pointwise confidence intervals.}
    \label{fig:qq}
\end{figure}

The Gaussian Q-Q plots show systematic departures in both tails for both LSTM models, with the left tail exhibiting larger departures than the right.
The Student-$t$ Q-Q plots show substantially improved tail agreement, confirming that the excess kurtosis in $z_t$ is better described by a heavy-tailed distribution.
I omit the corresponding P-P plots, as they are less sensitive to the tail departures that are the primary diagnostic interest here (Section~18.3.2).

\subsection{Interval coverage}

Table~\ref{tab:coverage} reports the empirical coverage at five nominal levels under both the Gaussian and Student-$t$ predictive laws.

\begin{table}[H]
\centering
\caption{Interval coverage: empirical vs nominal under Gaussian and Student-$t$ predictive laws.}
\label{tab:coverage}
\begin{ruledtabular}
\begin{tabular}{lcccccc}
& \multicolumn{2}{c}{\textbf{NLL LSTM}} & \multicolumn{2}{c}{\textbf{Composite LSTM}} & \multicolumn{2}{c}{\textbf{GARCH(1,1)}} \\
\hline
\textbf{Nominal} & Gauss & $t$ & Gauss & $t$ & Gauss & $t$ \\
\hline
50\% & 0.606 & 0.612 & 0.575 & 0.583 & 0.546 & 0.578 \\
80\% & 0.879 & 0.897 & 0.874 & 0.879 & 0.842 & 0.874 \\
90\% & 0.934 & 0.945 & 0.937 & 0.945 & 0.922 & 0.943 \\
95\% & 0.960 & 0.977 & 0.966 & 0.971 & 0.960 & 0.963 \\
99\% & 0.989 & 0.991 & 0.991 & 0.991 & 0.977 & 0.994 \\
\end{tabular}
\end{ruledtabular}
\end{table}

All three models show empirical coverage exceeding nominal coverage at every level, indicating that all models are overdispersed rather than overconfident.
To assess statistical significance, the 95\% binomial confidence interval for each nominal level is computed as $\alpha \pm 1.96\sqrt{\alpha(1-\alpha)/n}$ with $n = 348$.
At the 50\% level, the 95\% CI is $[0.447, 0.553]$: the NLL LSTM (0.606), composite LSTM (0.575), and GARCH (0.546) all significantly exceed nominal coverage, confirming systematic overdispersion.
At the 95\% level, the 95\% CI is $[0.927, 0.973]$: the NLL LSTM (0.960) falls within the interval, indicating that its 95\% intervals are statistically well-calibrated, while the composite LSTM (0.966) is borderline.
The Student-$t$ correction widens the intervals at extreme levels but narrows the effective intervals at central levels, generally improving the match to nominal coverage.

Figure~\ref{fig:coverage_curve} presents the calibration curve plotting empirical against nominal coverage for all three models.

\begin{figure}[H]
    \centering
    \includegraphics[width=0.8\textwidth]{../notebooks/daily/return_prediction/figures/figure6_coverage_calibration_curve.pdf}
    \caption{Calibration curve: empirical vs nominal coverage under the Gaussian predictive law.
    A perfectly calibrated model lies on the dashed 45-degree line.}
    \label{fig:coverage_curve}
\end{figure}

\subsection{GARCH(1,1) benchmark}

The fitted GARCH(1,1) parameters are $\hat{\omega} = 2.23 \times 10^{-5}$, $\hat{\alpha} = 0.146$, and $\hat{\beta} = 0.839$, yielding a persistence of $\hat{\alpha} + \hat{\beta} = 0.985$.
This high persistence is typical for daily financial returns and indicates that volatility shocks are long-lived.
The average predicted $\sigma_t$ is 0.0238, close to the empirical return standard deviation of 0.0230.

Despite having the lowest KS statistic ($D_n = 0.043$, $p = 0.534$), GARCH produces the highest AD statistic ($A^2 = 2.081$), the highest excess kurtosis in $z_t$ (11.64), and the most extreme skewness ($-1.40$).
This indicates that GARCH tracks the volatility level well on most days but fails catastrophically on extreme days --- exactly the pattern one would expect from a model that assumes Gaussian innovations.
The correlation between $\sigma_t^{\text{GARCH}}$ and $|r_t|$ is only 0.060, substantially lower than for the LSTM models, indicating that GARCH's one-step-ahead forecasts respond to yesterday's shock rather than anticipating today's volatility.

\subsection{Out-of-sample likelihood comparison}

Table~\ref{tab:nll} reports the out-of-sample NLL per observation and the AIC under both distributional assumptions.

\begin{table}[H]
\centering
\caption{Out-of-sample NLL per observation and AIC. Lower NLL indicates better fit.}
\label{tab:nll}
\begin{ruledtabular}
\begin{tabular}{lcccc}
\textbf{Model} & \textbf{Gauss NLL} & \textbf{Student-$t$ NLL} & \textbf{AIC (Gauss)} & \textbf{AIC ($t$)} \\
\hline
NLL LSTM       & $-2.332$ & $-2.349$ & $-1623.0$ & $-1632.8$ \\
Composite LSTM & $-2.195$ & $-2.200$ & $-1527.9$ & $-1529.0$ \\
GARCH(1,1)     & $-2.309$ & $-2.367$ & $-1607.3$ & $-1645.4$ \\
\end{tabular}
\end{ruledtabular}
\end{table}

The NLL LSTM achieves the best Gaussian out-of-sample likelihood ($-2.332$), outperforming GARCH ($-2.309$) and the composite LSTM ($-2.195$).
The Student-$t$ specification improves the NLL for all three models, with the largest improvement for GARCH ($-2.309 \to -2.367$), followed by the NLL LSTM ($-2.332 \to -2.349$) and the composite LSTM ($-2.195 \to -2.200$).
The AIC confirms this ranking: under the Gaussian specification, the NLL LSTM achieves the lowest AIC ($-1623.0$), followed by GARCH ($-1607.3$) and the composite LSTM ($-1527.9$).
Under the Student-$t$ specification, GARCH achieves the best AIC ($-1645.4$), reflecting the large kurtosis reduction from adopting heavy tails, followed by the NLL LSTM ($-1632.8$) and the composite LSTM ($-1529.0$).
The composite LSTM has the worst NLL and AIC under both specifications, indicating that the composite loss produces less accurate predictive distributions despite yielding better directional accuracy.

However, raw NLL differences do not establish statistical significance.
A Diebold--Mariano test \citep{diebold1995comparing} on the daily log-score differentials reveals that the NLL LSTM significantly outperforms the composite LSTM ($\bar{d} = +0.137$ nats/day, $t = 4.65$, $p < 0.001$), confirming that the composite loss degrades probabilistic quality in a statistically robust sense.
In contrast, the NLL LSTM does \emph{not} significantly outperform GARCH ($\bar{d} = +0.023$ nats/day, $t = 0.43$, $p = 0.67$).
This is a striking result: despite conditioning on 51 features through a deep recurrent architecture, the NLL LSTM's predictive density is statistically indistinguishable from a three-parameter GARCH model that uses only lagged returns.
The GARCH-versus-composite comparison is borderline ($p = 0.087$), suggesting that even GARCH tends to outperform the composite LSTM on probabilistic grounds.

To quantify the practical magnitude of these differences, the mean daily log-score gap between the NLL LSTM and the composite model is $0.197$ bits per day.
Over the 348-day test period, this corresponds to a cumulative log-likelihood ratio of $95.1$, which is overwhelmingly significant.
For the NLL LSTM versus GARCH, the daily gap is only $0.032$ bits per day --- a negligible difference relative to the daily volatility of the log-score differential.

\subsection{Student-$t$ extension}

Table~\ref{tab:student_t} reports the fitted degrees-of-freedom parameter and the KS test results under the Student-$t$ predictive law.

\begin{table}[H]
\centering
\caption{Student-$t$ extension: fitted degrees of freedom and KS test against $t_{\hat{\nu}}(0,1)$.}
\label{tab:student_t}
\begin{ruledtabular}
\begin{tabular}{lccc}
\textbf{Model} & $\hat{\nu}$ & $D_n$ \textbf{(KS)} & $p$\textbf{-value (KS)} \\
\hline
NLL LSTM       & 16.15 & 0.088 & 0.009 \\
Composite LSTM & 29.07 & 0.066 & 0.090 \\
GARCH(1,1)     & 10.62 & 0.051 & 0.313 \\
\end{tabular}
\end{ruledtabular}
\end{table}

The fitted $\hat{\nu}$ values confirm that the standardised residuals have heavier tails than a Gaussian: the NLL LSTM residuals require $\hat{\nu} \approx 16$ degrees of freedom, while the GARCH residuals require $\hat{\nu} \approx 11$.
The composite LSTM has the highest $\hat{\nu}$ ($\approx 29$), meaning its residuals are closest to Gaussian --- but this is because the composite loss inflates $\sigma_t$, which absorbs some of the tail mass and reduces the apparent kurtosis of the standardised residuals.

Figure~\ref{fig:t_qq} presents the Q-Q plots of $z_t$ against the fitted Student-$t$ distributions for both LSTM models.

\begin{figure}[H]
    \centering
    \includegraphics[width=\textwidth]{../notebooks/daily/return_prediction/figures/figure7_student_t_qq_plots.pdf}
    \caption{Q-Q plots of $z_t$ against the fitted Student-$t$ distribution: NLL LSTM (left, $\hat{\nu} = 16.15$) and composite LSTM (right, $\hat{\nu} = 29.07$).}
    \label{fig:t_qq}
\end{figure}

\subsection{Rolling calibration}

Figure~\ref{fig:rolling_cov} presents the rolling 60-day empirical coverage at the 95\% nominal level for all three models.
Figure~\ref{fig:rolling_corr} presents the rolling 60-day correlation between $\sigma_t$ and $|r_t|$.

\begin{figure}[H]
    \centering
    \includegraphics[width=\textwidth]{../notebooks/daily/return_prediction/figures/figure8_rolling_95_coverage.pdf}
    \caption{Rolling 60-day empirical coverage at the 95\% nominal level.
    The dashed line indicates the 0.95 reference.}
    \label{fig:rolling_cov}
\end{figure}

\begin{figure}[H]
    \centering
    \includegraphics[width=\textwidth]{../notebooks/daily/return_prediction/figures/figure9_rolling_sigma_abs_return_corr.pdf}
    \caption{Rolling 60-day correlation between predicted $\sigma_t$ and absolute returns $|r_t|$.
    Higher values indicate better real-time volatility tracking.}
    \label{fig:rolling_corr}
\end{figure}

Figure~\ref{fig:conditional} presents the conditional calibration analysis, splitting the test set into high-volatility ($n = 145$) and low-volatility ($n = 144$) subsets at the median rolling 60-day realised volatility (0.0227).
The results reveal that calibration quality is strongly regime-dependent.
In the low-volatility regime, all three models produce residuals that are approximately Gaussian: the excess kurtosis drops to $0.69$ for the NLL LSTM and $0.07$ for the composite LSTM, compared to $7.76$ and $5.19$ in the high-volatility regime.
The KS test fails to reject the Gaussian null for any model in either regime at the 5\% level, but the high-volatility kurtosis values confirm that the aggregate AD rejections are driven primarily by extreme days during volatile periods.
The NLL per observation is substantially worse in the high-volatility regime ($-2.217$ for the NLL LSTM) than in the low-volatility regime ($-2.464$), indicating that the LSTM's learned $\sigma_t$ does not fully capture the magnitude of volatility spikes.

This conditional analysis addresses a limitation of the unconditional tests: a model can appear well-calibrated on average while being systematically overconfident during the periods that matter most for risk management.

\begin{figure}[H]
    \centering
    \includegraphics[width=\textwidth]{../notebooks/daily/return_prediction/figures/figure10_conditional_calibration.pdf}
    \caption{Conditional calibration: standardised residual distributions in high-volatility (top) and low-volatility (bottom) regimes, split at the median rolling 60-day realised volatility.
    Annotations show the KS $p$-value and excess kurtosis for each model--regime combination.}
    \label{fig:conditional}
\end{figure}

The rolling coverage plot reveals that miscalibration is not constant over time but varies with the volatility regime.
The NLL LSTM's rolling 95\% coverage dips below the nominal level around mid-2025, coinciding with a period of elevated volatility in the WTI market.
This episode illustrates a known limitation of the LSTM's learned $\sigma_t$: because the model was trained on data ending in December 2022, it has never seen the specific volatility patterns of the 2025 test period and may underestimate uncertainty during novel regimes.
GARCH, by contrast, maintains more stable rolling coverage because its recursive structure continuously adapts to recent shocks via the persistence equation $\sigma_t^2 = \omega + \alpha r_{t-1}^2 + \beta \sigma_{t-1}^2$.

The rolling correlation plot shows that the NLL LSTM exhibits modest positive correlation with $|r_t|$ during most of the test period, confirming that it captures some contemporaneous volatility signal.
The composite LSTM's rolling correlation is predominantly negative, reinforcing the finding from Section~4.6 that the composite loss inverts the volatility-tracking relationship.
This negative correlation means that the composite LSTM's $\sigma_t$ is anti-correlated with realised volatility --- it widens intervals when the market is calm and narrows them when the market is volatile, precisely the opposite of what a risk manager would want.
GARCH maintains a stable but low positive correlation, reflecting its purely backward-looking volatility estimate.

\subsection{Christoffersen conditional coverage test}

The unconditional 95\% coverage test (Table~\ref{tab:coverage}) shows that all three models pass.
However, the Christoffersen independence test reveals a subtler failure.
For the NLL LSTM, the conditional probability of a violation following a previous violation is $\hat{\pi}_{11} = 0.214$, compared to $\hat{\pi}_{01} = 0.033$ for violations following no violation --- a sixfold increase.
The likelihood ratio statistic is $\text{LR}_{\text{ind}} = 6.11$ ($p = 0.013$), rejecting the independence null at the 5\% level.
In plain terms, when the NLL LSTM's 95\% interval fails on one day, it is substantially more likely to fail the next day as well.
This clustering indicates that $\sigma_t$ does not adapt fast enough to sudden volatility changes --- exactly what one would expect from a model that was trained on data ending in December 2022 and never updated.

The composite LSTM ($\hat{\pi}_{11} = 0.167$, $p = 0.058$) and GARCH ($\hat{\pi}_{11} = 0.143$, $p = 0.113$) both pass the independence test, though the composite is borderline.
GARCH's passing is noteworthy: its recursive structure allows it to adapt to yesterday's shock, preventing the violation clustering that affects the fixed-weight LSTM.

This finding adds an important nuance to the calibration picture: the NLL LSTM passes the unconditional coverage test but fails the conditional test, meaning its intervals are correct on average but not in the right places.
For risk management applications, where consecutive breaches compound losses, this distinction is critical.

\subsection{Position-sizing backtest}

To connect the calibration findings to the group's proposed future work on $\sigma_t$-based position sizing, I conduct a simple inverse-volatility backtest.
The position size on each day is $w_t = \sigma_{\text{target}} / \sigma_t$, where $\sigma_{\text{target}} = 0.02$ (daily), capped at $5\times$ leverage.
If $\sigma_t$ is perfectly calibrated, the realised portfolio volatility should equal the target.

The \emph{volatility ratio} (realised / target) provides a direct measure of calibration quality for sizing: the NLL LSTM achieves $0.931$, GARCH achieves $1.037$, and the composite LSTM achieves $0.826$.
The NLL LSTM and GARCH both produce portfolios with realised volatility close to the target, confirming that their $\sigma_t$ estimates are usable for risk budgeting.
The composite LSTM's ratio of $0.826$ indicates that its inflated $\sigma_t$ leads to systematically undersized positions and lower portfolio volatility than intended --- the model is too conservative by approximately 17\%.

\begin{figure}[H]
    \centering
    \includegraphics[width=\textwidth]{../notebooks/daily/return_prediction/figures/figure12_position_sizing_backtest.pdf}
    \caption{Inverse-volatility position sizing backtest.
    Left: cumulative portfolio returns.
    Right: rolling 60-day annualised portfolio volatility against the target (dashed line).
    The NLL LSTM and GARCH track the target closely; the composite LSTM consistently undershoots.}
    \label{fig:position_sizing}
\end{figure}

\subsection{Calibration scorecard}

Table~\ref{tab:scorecard} consolidates the pass/fail outcomes of all calibration tests at the 5\% significance level.
The scorecard reveals a clear pattern: GARCH and the composite LSTM pass more distributional tests than the NLL LSTM, but for opposite reasons.
GARCH passes the KS test because the KS statistic is insensitive to the tail departures that dominate its residuals.
The composite LSTM passes the KS test because its inflated $\sigma_t$ mechanically absorbs extreme residuals.
In both cases, the distributional test outcomes are misleading when interpreted in isolation --- they must be read alongside the NLL and Diebold--Mariano results, which show that the NLL LSTM produces the best (or statistically equivalent) predictive density.

\begin{table}[H]
\centering
\caption{Calibration scorecard: pass/fail at 5\% significance.}
\label{tab:scorecard}
\begin{ruledtabular}
\begin{tabular}{lccc}
\textbf{Test} & \textbf{NLL LSTM} & \textbf{Composite LSTM} & \textbf{GARCH(1,1)} \\
\hline
KS (Gaussian)         & Fail & Pass & Pass \\
AD (Gaussian)         & Fail & Fail & Fail \\
Ljung--Box $z_t$      & Pass & Pass & Pass \\
Ljung--Box $|z_t|$    & Pass & Pass & Pass \\
95\% coverage (Gauss) & Pass & Pass & Pass \\
KS (Student-$t$)      & Fail & Pass & Pass \\
95\% coverage ($t$)   & Fail & Pass & Pass \\
\hline
\textbf{Pass count}   & 3/7  & 6/7  & 6/7  \\
\end{tabular}
\end{ruledtabular}
\end{table}

The scorecard reveals a paradox that is central to this study: the model with the best probabilistic forecast quality (NLL LSTM, by the Diebold--Mariano test and out-of-sample NLL) fails the most distributional tests, while the model with the worst forecast quality (composite LSTM, by NLL) passes the most.
This paradox underscores the importance of using proper scoring rules alongside distributional tests --- goodness-of-fit tests alone can be misleading when models differ in their level of overdispersion.

%% =========================================================================
\section{Discussion}
\label{sec:discussion}
%% =========================================================================

\subsection{Assessment of hypotheses}

\textbf{H1: NLL LSTM residuals are approximately $\mathcal{N}(0,1)$.}
I find that this hypothesis receives partial support.
The NLL LSTM's standardised residuals have near-zero mean ($-0.070$), unit-range standard deviation ($0.938$), and no significant autocorrelation in either $z_t$ or $|z_t|$ (Ljung--Box $p = 0.48$ and $p = 0.70$ respectively).
The KS test rejects the Gaussian null ($p = 0.015$), and the AD test rejects more strongly ($A^2 = 1.456$).
The 95\% coverage is 96.0\%, close to nominal.
The group's informal claim that the NLL LSTM's $\sigma_t$ is ``calibrated'' is therefore partially supported at the level of first- and second-order properties, but the formal distributional tests reject Gaussianity due to excess kurtosis ($5.01$) and negative skewness ($-0.89$).

\textbf{H2: Calibration deteriorates under the composite loss.}
This hypothesis is supported, but with a subtlety I did not anticipate.
The composite LSTM's out-of-sample NLL ($-2.195$) is substantially worse than the NLL LSTM's ($-2.332$), confirming that the composite loss degrades the probabilistic quality of the forecasts.
The composite LSTM also shows borderline significant autocorrelation in $z_t$ ($p = 0.09$) and a standard deviation further from unity ($0.904$ versus $0.938$).
However, the KS test does not reject the Gaussian null for the composite LSTM ($p = 0.123$), and the AD statistic is lower ($1.160$ versus $1.456$).
This apparent paradox is explained by the fact that the composite loss inflates $\sigma_t$ (mean $0.031$ versus $0.025$ for the NLL LSTM), which absorbs tail exceedances and artificially reduces the kurtosis of the standardised residuals.
The composite LSTM thus appears ``better calibrated'' by some distributional metrics not because it is producing more accurate uncertainty estimates, but because it is producing wider intervals that mechanically reduce the apparent non-Gaussianity of the residuals.

This interpretation is confirmed by the out-of-sample NLL: the NLL LSTM assigns higher probability to the observed returns, indicating that its predictive distributions are tighter and more accurate.
The composite loss achieves better directional accuracy at the cost of worse probabilistic calibration, exactly as proper scoring rule theory predicts when auxiliary objectives are added to the likelihood.

\textbf{H3: A Student-$t$ predictive distribution improves calibration.}
This hypothesis is clearly supported.
The Student-$t$ extension improves the Q-Q tail agreement for both LSTM models (Figure~\ref{fig:t_qq}), improves the out-of-sample NLL ($-2.332 \to -2.349$ for the NLL LSTM), and improves coverage at extreme quantiles (Table~\ref{tab:coverage}).
The fitted $\hat{\nu}$ values (16 for NLL LSTM, 29 for composite, 11 for GARCH) confirm that standardised residuals exhibit persistent fat tails across all models.
This provides empirical support for the group's own future-work suggestion to replace the Gaussian likelihood with a heavier-tailed distribution.

\subsection{Comparison with GARCH(1,1)}

GARCH presents an instructive comparison that echoes the finding of \citet{hansen2005forecast} that simple GARCH(1,1) specifications are remarkably hard to beat.
Despite being a far simpler model from 1986 with only three parameters, it achieves a KS $p$-value of 0.534 (the highest of all models), the closest 95\% coverage to nominal (0.960), and no significant residual autocorrelation.
However, its AD statistic is the highest ($2.081$), and its residual kurtosis is extreme ($11.64$), indicating that while GARCH captures the volatility level well on normal days, it fails on extreme days.
The correlation between predicted $\sigma_t$ and same-day absolute returns $|r_t|$ --- a measure of real-time volatility tracking --- is $0.065$ for the NLL LSTM, $-0.043$ for the composite LSTM, and $0.060$ for GARCH.
The NLL LSTM's modest positive correlation indicates it captures some contemporaneous volatility signal from the 51-feature input set.
The composite LSTM's \emph{negative} correlation is notable: it implies that the composite loss has inverted the volatility-tracking relationship, increasing $\sigma_t$ when actual volatility is low and decreasing it when volatility is high.
This is a direct consequence of the variance-collapse penalty, which penalises low $\text{Var}(\mu_t)$ and indirectly distorts $\sigma_t$.
GARCH's correlation ($0.060$) reflects the one-step lag inherent in the recursion: today's forecast depends on yesterday's shock.

For uncertainty quantification specifically, the NLL LSTM achieves a better raw NLL ($-2.332$ versus $-2.309$), but the Diebold--Mariano test shows that this difference is not statistically significant ($p = 0.67$).
This result is consistent with the broader finding of \citet{hansen2005forecast} that GARCH(1,1) is remarkably hard to beat for volatility forecasting.
One explanation is that the LSTM was trained on data ending in December 2022 and never updated, while GARCH's recursive structure continuously adapts to new information.
A walk-forward retraining protocol might reveal a larger LSTM advantage, but the current evidence does not support the conclusion that the LSTM provides a statistically significant improvement in probabilistic forecast quality over GARCH.

\subsection{Implications for position sizing}

The group project proposed using $\sigma_t$ for position sizing.
The inverse-volatility backtest (Section~4.11) provides direct evidence on this question.
The NLL LSTM achieves a volatility ratio of $0.931$ --- meaning realised portfolio volatility is within 7\% of the target --- which I consider adequate for practical use.
However, the Christoffersen test (Section~4.10) reveals that the NLL LSTM's violations are clustered, meaning consecutive breaches are more likely during volatile periods.
In my view, this is the most important practical finding: the NLL LSTM's $\sigma_t$ is a reasonable sizing input on average, but a risk manager should be aware that it underreacts to sudden volatility spikes, and a GARCH overlay or regime-switching adjustment would help address this.
The composite LSTM's $\sigma_t$ should not be used for sizing: it overestimates risk by 17\% and is anti-correlated with realised volatility.

\subsection{Limitations}

Several limitations qualify the conclusions beyond the test-set reuse noted in Section~\ref{sec:methodology}.
First, the analysis covers a single asset (WTI crude oil) during a specific market regime; the findings may not generalise to other commodities, equities, or time periods with different volatility dynamics.
Second, the conditional calibration analysis uses an in-sample median split rather than a pre-specified volatility threshold, which introduces a mild look-ahead bias in the regime classification.
Third, the Diebold--Mariano test assumes that the log-score differentials are stationary, an assumption that may be violated during regime transitions.
Fourth, the Student-$t$ extension fits $\hat{\nu}$ to the full test set, whereas in a live deployment, $\hat{\nu}$ would need to be estimated on a rolling basis, introducing additional estimation noise.
Finally, the calibration scorecard paradox (Section~4.10) highlights a broader conceptual issue: there is no single criterion that captures all dimensions of probabilistic forecast quality, and the choice of evaluation metric can reverse the apparent ranking of models.

%% =========================================================================
\section{Conclusion and Future Work}
\label{sec:conclusion}
%% =========================================================================

This report applied a formal calibration battery to the probabilistic outputs of two LSTM variants and a GARCH(1,1) benchmark, testing the group project's informal claim that the LSTM's learned $\sigma_t$ is calibrated.

The main findings are as follows.
The NLL LSTM produces partially calibrated uncertainty estimates: the standardised residuals have approximately correct first and second moments, no significant autocorrelation, and close-to-nominal coverage at central quantiles.
However, the Gaussian distributional assumption is formally rejected due to excess kurtosis and negative skewness, with the conditional calibration analysis showing that this rejection is driven primarily by the high-volatility regime.
The composite loss significantly degrades probabilistic calibration ($p < 0.001$ by Diebold--Mariano test), producing worse out-of-sample likelihood despite better directional accuracy, consistent with proper scoring rule theory.
Perhaps the most sobering finding is that the NLL LSTM does not significantly outperform the three-parameter GARCH model ($p = 0.67$) on probabilistic forecast quality, despite its far greater complexity and access to 51 features.
A Student-$t$ extension consistently improves calibration diagnostics for all models, confirming the group's suggestion on an empirical basis.

Future work could pursue four directions.
First, training the LSTM directly under a Student-$t$ negative log-likelihood, rather than post-hoc fitting the degrees-of-freedom parameter to Gaussian residuals, would yield an end-to-end heavy-tailed probabilistic model.
Second, a walk-forward retraining protocol would address the stationarity concern: the LSTM was trained on data ending in December 2022 and never updated, while GARCH's recursive structure continuously adapts, potentially explaining their comparable performance.
Third, a formal economic evaluation using $\sigma_t$-based position sizing would quantify whether the calibration improvements documented here translate into improved risk-adjusted returns.
Fourth, extending the analysis to multiple assets and time periods would establish whether these findings generalise beyond WTI crude oil.

\bibliography{references}

\appendix
\section{Appendixes}
This appendix is reserved for supplementary material.

\end{document}
